\section{2026-08-22 — La encuesta y una brecha que no se esperaba}

Cerró el campo de la encuesta con 140 respuestas, de las cuales 107 satisfacen los tres criterios de segmento. El número total supera el mínimo de 120 que pidió la cátedra, aunque conviene tener presente que la base de análisis efectiva es la de 107 y que las tablas del anexo lo indican en cada caso. La distribución se hizo por vía digital y presencial, lo que explica que una parte de las respuestas se cargara de forma manual a la planilla.

El análisis se automatizó en un guion propio que lee la planilla y emite tanto las tablas del wiki como las figuras del documento en formato de \emph{pgfplots}. La decisión no fue estética. La cátedra pidió expresamente que no se peguen capturas del formulario, y un guion que reconstruye las figuras a partir de la fuente evita el riesgo de que los porcentajes del texto y los de los gráficos se desincronicen ante cualquier corrección posterior. La contrapartida es que las figuras no deben editarse a mano, cosa que se dejó anotada en la cabecera de cada archivo generado.

Los cuatro supuestos que motivaron la encuesta quedaron validados, pero el hallazgo que importa no era ninguno de ellos. La intención declarada de usar la extensión promedia 4,06 sobre 5 mientras que la confianza en el puntaje automático promedia 3,57. Las personas encuestadas quieren la herramienta bastante más de lo que le creen al número que produciría. El cruce con la capacidad autopercibida de detección refuerza la lectura: la confianza en el puntaje desciende de forma monótona a medida que sube la capacidad que la persona se atribuye.

La consecuencia sobre el producto es concreta y obliga a revisar el énfasis de los capítulos ya escritos. Hasta ahora la evidencia y la explicación aparecían como una prestación más de la interfaz, ubicada al mismo nivel que las demás. La encuesta indica que son la condición de uso: un puntaje entregado sin la evidencia que lo sostiene no resulta utilizable para el segmento relevado. Conviene que la defensa apoye ahí y no en la métrica del clasificador.

El segundo hallazgo con consecuencias de diseño es la centralidad de la neutralidad política. Encabeza los atributos de confianza con un 84,1 por ciento, veinte puntos por encima del segundo, y reaparece del lado de las preocupaciones junto con el temor a que se censuren opiniones o contenido satírico. Leído en conjunto con el 66,4 por ciento que teme al falso positivo, el mensaje es que el riesgo percibido no es la incapacidad técnica sino la toma de partido. Eso ordena la prioridad del umbral de decisión hacia la precisión por encima de la exhaustividad, criterio que ya estaba enunciado en la sección de validación pero sin fundamento empírico propio.

Queda pendiente el instrumento cualitativo. Las dos entrevistas todavía no se realizaron y, en consecuencia, la persona secundaria y el anexo de transcripciones siguen sin contenido. El anexo correspondiente permanece comentado en el archivo de anexos.

\bigskip
